%%Tair_soni_ha
Etant donn� ces 3 formules thermodynamiques :
\begin{align*}
T &= (1 - 0.518.q).T_{soni}\\
q &= r/(1+r) \\
r &= H_a.R_a.T/P 
\end{align*}

o� $q$ est l'humidit� sp�cifique, $r$ le rapport de m�lange en vapeur d'eau,
$R_a$ est la constante sp�cifique de l'air sec.\\

On en d�duit que la temp�rature de l'air $T$ est une racine du polyn\^ome du 
second degr� suivant :
\[
T^2 + \left [ \frac{P}{H_a.R_a} - (1 - 0.518)\times T_{soni} \right ].T
    - \frac{P}{H_a.R_a}.T_{soni} = 0  
\]

